\section{Discussion}\label{sec:discussion}

\subsection{Three Barriers, Only the First of Which a Prompt Can Remove}
\label{sec:barriers}

The two prompt conditions differ by only one variable, and the effect of that variable
cleanly splits between languages. For Python, adding the target specification cut the number
of suites that failed to import or collect from 17 to 4 --- 14 targets rescued and 1 newly
broken. For Java, it made virtually no difference: 38 failed to compile with the baseline
prompt, and 38 with the target-specified one, with 35 functions overlapping between those
sets. Three functions moved in either direction, which is noise, not an effect.

The difference in effect between languages is not a peculiarity of their tools --- it is
encoded in the block itself. The block specifies the identity of the target: its qualified
name, declaring class or module, signature, and a constructor call for the receiver, if any.
For Python, that is enough to remove barriers to invocation --- at the stage these suites
failed, the failure was always importing a module from an incorrect path, calling a
non-existent function, or attempting to invoke a method on a class that does not support it.
Give Python the correct name and location, and it will resolve the symbol successfully. Java
performs these checks at compile time, so a correct name is still insufficient --- the model
must synthesise type-correct uses of the symbol within the suite, including receiver arguments
of the correct arity and parameter types, return types, thrown exceptions, generics, and so
on. The target block does not include any of that information.

This suggests a three-level reading of the results in terms of what the generated test must
do to observe the target. We offer this as a hypothesis for future work:

\begin{enumerate}
  \item \textbf{Addressing} --- calling the right thing at all. The target specification
    removed this barrier, almost entirely in Python.
  \item \textbf{Typing} --- calling it in a way the compiler accepts. Only relevant to Java.
    The target specification does not meaningfully help, because the difficulty is not with
    the target itself, but the surrounding API.
  \item \textbf{State} --- constructing a receiver in a configuration where the behaviour of
    interest is visible. Neither prompt condition addresses this.
\end{enumerate}

The reason to take this seriously is the evidence in support of the third item. For Python,
the number of suites that imported successfully but had no passing test barely moved between
the two prompt conditions: 13 under the baseline prompt, 14 under the target-specified one,
with only 7 functions appearing in both sets. Removing the addressing barrier did not shrink
the state barrier; it \emph{exposed} it. The population changed while the size did not.

The observation that the barrier is not with assertion writing but with receiver construction
is not new --- it is the reason search-based test generation struggles with object-oriented
code~\cite{fraser2014largescale,shamshiri2015realfaults}. The point of this finding is that
the same wall exists for language models, and that providing the constructor is not enough to
get past it --- our snippets were run and checked to create the receiver, and the state
barrier persisted. Building an instance is not the same as building an instance in a state
where some behaviour is triggered.

\subsection{Reaching a Target and Testing It Are Different Achievements}

The clearest finding of this study is the difference between two tiers that would have
appeared identical in a single-metric study. On Python, the target specification increased the
proportion of suites that reached the target from 29/60 to 42/60 ($p = 0.003$, $r_{rb} =
+0.61$), but its effect on the number that could detect an injected fault was not significant,
16/60 to 22/60 ($p = 0.33$).

A study that used only coverage as a metric would have concluded that the enriched prompt is
an improvement across the board, while also being wrong about what aspect of testing it
improved. The suites arrived at the right function more often, but tested it no more often.
Coverage measures arrival and mutation score measures discrimination --- only the latter is a
metric of testing ability in any meaningful sense, and only the former was improved by the
prompt change. That coverage is misleading about test-suite quality is a well-known
finding~\cite{inozemtseva2014coverage}; what this experiment adds is a demonstration that a
specific, plausible change to the prompt moves the two metrics by different amounts, and in a
way that inverts the conclusion drawn from them.

This is also where the pre-registered hypothesis H-B begins to fall apart. We predicted, with
no knowledge of the results, that the enriched prompt would be no better than the baseline.
That prediction holds at T4 and fails at T3. We report this as a partially refuted hypothesis,
rather than an overturned one, because the direction of effect at T3 is the one we predicted
against.

\subsection{Neither Family Wins Outright}

The comparison with existing generators also fails to reveal a clear winner, and the
languages rank oppositely.

On Python, both prompt conditions substantially outperformed Pynguin: 16/60 and 22/60 suites
detecting a fault compared to 2/60, $p \le 0.003$ in all pairs, with rank-biserial
correlations ranging from $-0.81$ to $-0.96$. The reason for the gap was apparent in the
numbers of failed suites: Pynguin failed to generate any suite for 39/60 targets in its
default configuration. The documented workaround for the reported serialisation crash
(disabling the master--worker split) reduced this to 35, a net movement of four that conceals
considerable churn: 17 targets gained a suite and 13 lost one, and 14 of the 17 new suites
failed at collection, leaving T4 at 2/60. We report this result because it is the sort of
finding that disappears when a study only reports the better-performing configuration. In this
case, neither configuration worked.

On Java, the count ordering was reversed: the search-based tools found faults on more targets
than the model in all pairs (Section~\ref{sec:rqc}). The reversal was firmer for EvoSuite, for
which the paired test reached significance at the whole-suite gate ($p = 0.028$), than for
Randoop, for which no gate did; we interpret this as a consistent count advantage for the
mature tools, rather than a uniformly significant one. Java is also where these tools were
developed and where they have been iterated upon for the last decade, and where static typing
gives them the information needed to construct calls --- the information the language model was
demonstrated to lack in Section~\ref{sec:barriers}. Python removes that information, and the
ordering flips.

The conclusion is that neither family of tools is demonstrably better overall, but that they
fail in ways that are specific to languages. A study that examined only Java would have
confidently reported the superiority of search-based testing; a study that examined only
Python would have confidently reported the opposite.

\subsection{What This Means for Evaluating Test Generators}

Three practices follow from what went wrong rather than from what worked.

\paragraph{Report the sampling funnel.}
Only 26\,\% of medium-complexity functions in these ten projects are both public and
unambiguously named. A benchmark that samples on complexity alone will therefore spend most of
its mass on targets that no external generator can address, and every tool evaluated on it is
charged for failures that belong to the sampling rather than to the tool. The funnel costs
nothing to report, and it bounds how much of a measured failure rate can be attributed to the
tool at all.

\paragraph{Rest conclusions on a tier that cannot be satisfied vacuously.}
Both compilation success and the green-test criterion were defeated in this corpus.
Compilation is defeated by reflection, which turns names into unchecked strings. The
green-test criterion is defeated by a suite that makes a single assertion-free call to the
wrong function: it compiles, it collects, and it passes, so the criterion records it as a
success while branch coverage inside the target's line range is zero. Neither case requires an
adversary; both appeared in ordinary tool output. Metrics should be interpreted with extra
caution in the presence of reflection, or similar techniques that let code bypass static
checks.

\paragraph{Assume the harness is broken until a failure proves otherwise.}
The failure mode that consumed the most time in this study was not a tool underperforming, but
a metric misreporting. Hardcoded constants, validity flags erroneously indicating success,
reflection removing the compiler's ability to check names, a version mismatch reporting no
tests while exiting successfully, a JDK mismatch doing the same, and a green-test gate
implemented differently in the two languages under one name appeared six times in total, four
of which were introduced by us while attempting to fix the other two. None of these issues
caused an exception to be thrown, and consequently, none of them were caught by the initial
harness design. The requirement for a harness is that it must have a test that fails --- if the
pipeline below the metric is entirely broken but the metric reports no change, that metric is
no longer useful. We enforce this requirement for all reported rates, and it is the cheapest
such check we know of.
